\scanpage{72}
Si d'autre part $\underline G$ est un schéma en groupes fini étale sur $S$, et si $X'$ est muni d'une structure de $\underline G_X$-torseur au dessus de $X^\ast$, alors le stabilisateur
% slanted
\margin{\uline{NB} On suppose que $\underline G$ opère sur $X'$, de façon que $X'^\ast=X'|X^\ast$ devienne un $\underline G_X$-torseur sur $X^\ast$}
dans $\underline G_{\deleted{\unsure{X|S}}}$ \deleted{\unsure{du}} de la section $t'$ est un sous-schéma en groupes \inserted{$\underline G_{t'}$} de $\underline G$, \deleted{\unsure{comme groupe|comme grou}} \inserted{qui} \deleted{\unsure{image dans|ment à|munit de}} \inserted{opère \inserted{fid.} sur} $T_{X'/S,t'}$ \uline{et} dont l'image dans $\deleted{T_{X'/S,t'}}\,\mathrm{Gm}_S\simeq\underline{\mathrm{Aut}}_{O_S}T_{X'/S,t'}$ est $\mu_{n\,S}$, d'où un monomorphisme canonique
\[
\mu_{n\,S}\hookrightarrow\underline G . \tag*{(3)}
\]

Par ailleurs, désignant par
\[
f^\circ:X^\ast\longrightarrow S
\]
le morphisme induit par $f$, le $\underline G_{X^\ast}$-torseur $X'^\ast$ définit une section
\[
\xi_{X'/S}\in H^0(S,R^1 f^\circ_\ast f^{\circ\ast}(\underline G)) . \tag*{(4)}
\]

Tout ceci posé, on a :

\uline{Proposition} \quad $(X',t')$ au dessus de $(X,\unsure{t|f})$, avec l'opération de $\underline G$ dessus, est défini à isom unique près par la connaissance de (3) (4) et de (2). \deleted{\unsure{(ces données|(ces derniers|(ce dernier}} De façon précise,
% archive sheet number: 71
